\subsection{A bordered-permanent induction criterion}

This is an exact conditional reduction. The finite tests described
below do not establish its hypothesis for every matrix size.

Let $E=\mathbb C^3$, and let $A(x)$ be an $m\times m$ matrix of
linear forms on $E$. Write $P(A)=\operatorname{per}(A)$ and
$G_{ij}(A)=\partial P/\partial A_{ij}$. For vectors $u(x),v(x)$
of linear forms, define
\[
 Q_A(u,v)=DP_A[uv^t]
          =\sum_{i,j}u_iv_jG_{ij}(A),\qquad
 R_A(u,v)=\tfrac12D^2P_A[uv^t,uv^t].
\]
Thus $P(A+tuv^t)=P(A)+tQ_A(u,v)+t^2R_A(u,v)+O(t^3)$.
The derivatives in $u,v$ below mean differentiation in independent
auxiliary coordinates, followed by substitution of the indicated
linear forms. Put
\begin{align*}
 a_i&=\partial_{u_i}Q_A=(Gv)_i,&
 b_j&=\partial_{v_j}Q_A=(G^tu)_j,\\
 t_i&=\partial_{u_i}R_A
          =D^2P_A[e_iv^t,uv^t],&
 t'_j&=\partial_{v_j}R_A
          =D^2P_A[ue_j^t,uv^t].
\end{align*}
All forms in
\begin{equation}\label{border:C}
 \mathcal C_A(u,v)
 =E^*P(A)
   +\sum_i E^*a_i+\sum_j E^*b_j
   +\operatorname{span}\{t_i,t'_j:1\leq i,j\leq m\}
 \ \subseteq\operatorname{Sym}^{m+1}E^*
\end{equation}
have the same degree. This space is generated by $8m+3$ displayed
columns. If at least one of $u,v$ is nonzero, the identity
\begin{equation}\label{border:scaling_relation}
 \sum_i u_i a_i-\sum_jv_jb_j=0
\end{equation}
is a nonzero relation on the corresponding $6m$ linear-multiple
columns. In particular $\dim\mathcal C_A(u,v)\leq8m+2$.

\paragraph{Conditional induction statement.}
If some ternary data $A,u,v$ satisfy
\begin{equation}\label{border:rank_hypothesis}
 \dim\mathcal C_A(u,v)=8m+2,
\end{equation}
then, with $c_m=\dim\mathcal P_{m,4}$,
\begin{equation}\label{border:increment}
 c_{m+1}\geq c_m+8m+2.
\end{equation}
Consequently, if $m\geq17$ and $c_m=4m^2-2m+2$, the same
equality holds at $m+1$. A proof of
\eqref{border:rank_hypothesis} for all $m\geq17$, together with
R-24, would prove the exact parameter dimension for all those sizes.

\begin{proof}
Adjoin a fourth variable $\ell$, so $V^*=E^*\oplus\mathbb C\ell$.
Extend $A,u,v$ to matrices and vectors of quaternary linear forms.
The condition \eqref{border:rank_hypothesis} on their restrictions
to $\ell=0$ is a nonempty Zariski open condition on these
quaternary parameters. The condition that the differential image
\[
 T_A=\{DP_A[U]:U\in\operatorname{Mat}_m(V^*)\}
\]
have dimension $c_m$ is also nonempty and Zariski open. The
parameter space is irreducible, so both conditions hold
simultaneously for some quaternary $A,u,v$. No maximal-rank
assumption is made on the initial ternary matrix itself.

Consider the single family of quaternary pencils
\[
 B_\varepsilon=
 \begin{pmatrix}A&\varepsilon u\\
                 \varepsilon v^t&\ell\end{pmatrix}.
\]
Every permutation uses either the bottom-right entry, or one
entry from each off-diagonal block. Hence exactly
\begin{equation}\label{border:exact_permanent}
 \operatorname{per}(B_\varepsilon)
       =\ell P(A)+\varepsilon^2Q_A(u,v).
\end{equation}
Differentiation in the bottom-right entry gives every form in
$V^*P(A)$. Differentiation in the upper-right and lower-left
entries, divided by $\varepsilon$, gives every form in
$\sum V^*a_i+\sum V^*b_j$, for $\varepsilon\ne0$.

The remaining $2m$ columns come from exact cancellations, not
from independently chosen deformations. Fix a constant vector
$h\in\mathbb C^m$. Vary the upper-right block by $\ell h$ and
the upper-left block by $-\varepsilon hv^t$, leaving all other
entries fixed. The differential of
\eqref{border:exact_permanent} in that direction is
\begin{align*}
 \varepsilon\ell\,DP_A[hv^t]
  -\varepsilon\ell\,DP_A[hv^t]
  -\varepsilon^3D_AQ_A(u,v)[hv^t]
   &=-\varepsilon^3D^2P_A[hv^t,uv^t]\\
   &=-\varepsilon^3\sum_i h_it_i.
\end{align*}
Every entry of this parameter variation is a linear form, as
required. Varying instead the lower-left block by $\ell h^t$
and the upper-left block by $-\varepsilon uh^t$ similarly gives
$-\varepsilon^3\sum_jh_jt'_j$. Division by
$-\varepsilon^3$ is permitted in the differential image for
every nonzero $\varepsilon$.

Choose $8m+2$ of the fixed polynomial columns in
\eqref{border:C} whose restrictions to $\ell=0$ are independent.
All belong to the differential image of this same
$B_\varepsilon$, for every $\varepsilon\ne0$.
Choose constant quaternary variations $U_1,\ldots,U_{c_m}$ such
that $DP_A[U_i]$ form a basis of $T_A$. Their derivatives at
$B_\varepsilon$ are
\[
 \ell DP_A[U_i]+\varepsilon^2D_AQ_A(u,v)[U_i].
\]
At $\varepsilon=0$ these become a basis of $\ell T_A$.
The combined $c_m+8m+2$ limit columns are independent: restriction
to $\ell=0$ first annihilates all coefficients of the selected
normal columns, and multiplication by the nonzero polynomial
$\ell$ is injective on $T_A$.

A coefficient minor of the combined columns is therefore a
polynomial in $\varepsilon$ with nonzero constant term. It is
nonzero at some $\varepsilon\ne0$, proving
\eqref{border:increment}. Finally the torus upper bound increases
by exactly
\[
 [4(m+1)^2-2(m+1)+2]-[4m^2-2m+2]=8m+2,
\]
so equality propagates whenever the hypothesis is available.
\end{proof}

\paragraph{Exact finite tests and the remaining gate.}
The source \texttt{check\_border\_ternary.cpp} constructs the
coefficient evaluations in \eqref{border:C} by subset recurrence
for $P(A+tuv^t)$ through order two and reverse differentiation.
For dense integer ternary data, the tests at $m=11$ and $m=12$
give ranks $90$ and $98$, respectively, the full $8m+2$.
For the tested cyclic ternary data the ranks are instead $71$
at $m=11$ and $107$ at $m=17$. These finite computations are
exploratory evidence about the criterion; they do not establish
the hypothesis at other sizes or classify the cyclic case.

